\begin{enumerate}[label=\thesubsection.\arabic*, ref=\thesubsection.\theenumi]

% 1. Used in Q1 (Dice roll PMF / Z-transform)
\item Probability Mass Function (PMF) and Z-Transform of Discrete Random Variables:
\begin{align}
M_Y(Z) = E[Z^{-Y}] = \sum_k P_Y(k) Z^{-k} \label{eq:CS12025-pmf_ztransform}
\end{align}

% 2. Used in Q1 (Dice roll triple product)
\item Z-Transform of Sum of Independent Random Variables $Y = \sum_{i=1}^n Y_i$:
\begin{align}
M_Y(Z) = \prod_{i=1}^n M_{Y_i}(Z) \label{eq:CS12025-ztransform_sum}
\end{align}

% 3. Used in Q2 & Q11 (Fake coin problem)
\item Bayes' Theorem:
\begin{align}
\Pr (E_1|A) = \frac{\Pr (E_1) \Pr (A|E_1)}{\sum_{i=1}^2 \Pr (E_i) \Pr (A|E_i)} \label{eq:CS12025-bayes_theorem}
\end{align}
\solution
Bayes' theorem relates the conditional probability of event $E_1$ given evidence $A$ to the prior probabilities $\Pr(E_1)$ and $\Pr(E_2)$ and the likelihoods $\Pr(A|E_1)$ and $\Pr(A|E_2)$:
\begin{align}
\Pr (E_1|A) &= \frac{\Pr(A \cap E_1)}{\Pr(A)} \\
            &= \frac{\Pr (E_1) \Pr (A|E_1)}{\Pr (E_1) \Pr (A|E_1) + \Pr (E_2) \Pr (A|E_2)} \\
            &= \frac{\Pr (E_1) \Pr (A|E_1)}{\sum_{i=1}^2 \Pr (E_i) \Pr (A|E_i)}
\end{align}

% 4. Used in Q3 (Function composition)
\item Composition of Functions and Surjectivity (Onto Property):
\begin{align}
(f \circ g)(x) = f(g(x)), \quad \forall x \in A \label{eq:CS12025-function_composition}
\end{align}

% 5. Used in Q4 (Spanning tree shortest path)
\item Shortest Path Inequality in Subgraphs and Spanning Trees:
If $T$ is a minimum spanning tree of graph $G$, then for any vertices $u, v \in V(G)$:
\begin{align}
d_1(u, v) \le d_2(u, v) \label{eq:CS12025-shortest_path_mst}
\end{align}
where $d_1(u,v)$ is the shortest path in $G$ and $d_2(u,v)$ is the shortest path in $T$.

% 6. Used in Q5 (Recurrence relation z-transform property)
\item Z-Transform Property of $n a^n u(n)$:
\begin{align}
\mathcal{Z}\{n a^n u(n)\} = \frac{a z^{-1}}{(1 - a z^{-1})^2} \label{eq:CS12025-ztransform_n_an}
\end{align}
\solution
From the foundational definition of the unilateral Z-transform, the Z-transform of the sequence $x(n) = n a^n u(n)$ is given by the infinite series summation:
\begin{align}
\mathcal{Z}\{n a^n u(n)\} &= \sum_{n=0}^{\infty} n a^n z^{-n} \\
                          &= \sum_{n=0}^{\infty} n (a z^{-1})^n
\end{align}

Let $x = a z^{-1}$. The summation can be expressed as:
\begin{align}
S = \sum_{n=0}^{\infty} n x^n
\end{align}

Taking the standard geometric series identity, which converges for $|x| < 1$:
\begin{align}
\sum_{n=0}^{\infty} x^n = \frac{1}{1 - x}
\end{align}

Differentiating both sides with respect to $x$:
\begin{align}
\frac{d}{dx} \left( \sum_{n=0}^{\infty} x^n \right) &= \frac{d}{dx} \left( \frac{1}{1 - x} \right) \\
\sum_{n=0}^{\infty} n x^{n-1} &= \frac{1}{(1 - x)^2}
\end{align}

Multiplying both sides by $x$ to match the power of $n x^n$:
\begin{align}
\sum_{n=0}^{\infty} n x^n &= \frac{x}{(1 - x)^2}
\end{align}

Finally, substituting back $x = a z^{-1}$ into the derived expression yields the closed-form Z-transform:
\begin{align}
\mathcal{Z}\{n a^n u(n)\} &= \frac{a z^{-1}}{(1 - a z^{-1})^2}
\end{align}

% 7. Used in Q6 (Linear system consistency)
\item System of Linear Equations Consistency Analysis:
For a linear system $\mathbf{A}\mathbf{x} = \mathbf{b}$, the existence and uniqueness of solutions are determined by rank conditions on $\mathbf{A}$ and the augmented matrix $[\mathbf{A} \mid \mathbf{b}]$:
\begin{align}
\begin{cases}
\text{Unique Solution:} & \text{rank}(\mathbf{A}) = \text{rank}([\mathbf{A} \mid \mathbf{b}]) = n \\
\text{Infinite Solutions:} & \text{rank}(\mathbf{A}) = \text{rank}([\mathbf{A} \mid \mathbf{b}]) < n \\
\text{No Solution:} & \text{rank}(\mathbf{A}) \neq \text{rank}([\mathbf{A} \mid \mathbf{b}])
\end{cases}
\label{eq:CS12025-linear_system_rank}
\end{align}

% 8. Used in Q9 (Ternary strings without consecutive duplicates)
\item Number of Strings over Alphabet Size $k$ without Adjacent Duplicate Symbols:
\begin{align}
N_{\text{no\_consec}} = k(k-1)^{n-1} \label{eq:CS12025-no_consec_strings}
\end{align}
\solution
For a string of length $n$ defined over an alphabet of size $k$:
\begin{itemize}
    \item The first character can be chosen in $k$ different ways.
    \item For each subsequent position from $2$ to $n$, the character must be distinct from the immediately preceding character, giving $(k-1)$ choices per position.
\end{itemize}
By the fundamental counting principle, the total number of such strings is:
\begin{align}
N_{\text{no\_consec}} = k \times \underbrace{(k-1) \times (k-1) \times \dots \times (k-1)}_{n-1 \text{ times}} = k(k-1)^{n-1}
\end{align}

% 9. Used in Q10 (Piecewise function boundary conditions)
\item Boundary Continuity and Differentiability Conditions for Piecewise Functions:
\begin{align}
\lim_{x \to x_0^-} f(x) &= \lim_{x \to x_0^+} f(x) = f(x_0) \label{eq:CS12025-piecewise_continuity} \\
\lim_{x \to x_0^-} f'(x) &= \lim_{x \to x_0^+} f'(x) \label{eq:CS12025-piecewise_differentiability}
\end{align}

% 10. Used in Q12 (Bubble sort swaps)
\item Total Inversions (Swaps) in a Reverse-Sorted Array of Size $n$:
\begin{align}
\text{Swaps} = \frac{n(n-1)}{2} \label{eq:CS12025-bubble_sort_swaps}
\end{align}
\solution
For an array of size $n$ sorted in strictly descending order, every adjacent pair satisfies $A[j] > A[j+1]$. The total number of swap operations is the sum of comparisons across all passes:
\begin{align}
\text{Swaps} &= \sum_{i=0}^{n-2} (n - i - 1) \\
             &= (n - 1) + (n - 2) + \dots + 2 + 1 \\
             &= \frac{n(n-1)}{2}
\end{align}

% 11. Used in Q14 (Eigenvalues of matrix powers)
\item Matrix Characteristic Equation and Eigenvalues of Powers:
The eigenvalues $\lambda$ of $\mathbf{A}$ satisfy $\det(\mathbf{A} - \lambda \mathbf{I}) = 0$. For any positive integer $k$:
\begin{align}
\mathbf{A}\mathbf{x} = \lambda \mathbf{x} \implies \mathbf{A}^k \mathbf{x} = \lambda^k \mathbf{x} \label{eq:CS12025-eigenvalue_power}
\end{align}
\solution
Given the eigenvalue definition $\mathbf{A}\mathbf{x} = \lambda \mathbf{x}$, multiplying by $\mathbf{A}$ successively $k$ times gives:
\begin{align}
\mathbf{A}^2 \mathbf{x} &= \mathbf{A}(\mathbf{A}\mathbf{x}) \\
                        &= \mathbf{A}(\lambda \mathbf{x}) \\
                        &= \lambda (\mathbf{A}\mathbf{x}) \\
                        &= \lambda^2 \mathbf{x}
\end{align}
Extending this result by mathematical induction for any positive integer $k$:
\begin{align}
\mathbf{A}^k \mathbf{x} = \lambda^k \mathbf{x}
\end{align}

% 12. Used in Q16 (BFS tree height)
\item Minimum Height of Breadth-First-Search (BFS) Tree given Graph Diameter $D$:
\begin{align}
h(T) \ge \left\lceil \frac{D}{2} \right\rceil \label{eq:CS12025-bfs_tree_height}
\end{align}

% 13. Used in Q18 (Basic block leaders)
\item Basic Block Leader Criteria in Compiler Intermediate Code:
\begin{align}
\text{Leaders} = \{\text{First Instruction}\} \cup \{\text{Targets of Jumps}\} \cup \{\text{Instructions Following Jumps}\} \label{eq:CS12025-basic_block_leaders}
\end{align}

% 14. Used in Q21 (PDF normalization)
\item Continuous Probability Density Function (PDF) Normalization:
\begin{align}
\int_{-\infty}^{\infty} P(x) \, dx = 1 \label{eq:CS12025-pdf_normalization}
\end{align}

% 15. Used in Q21 (CDF range probability)
\item Cumulative Distribution Function (CDF) and Range Probability:
\begin{align}
F_X(x) = \int_{-\infty}^x P(t) \, dt, \quad P(a \le X \le b) = F_X(b) - F_X(a) = \int_a^b P(x) \, dx \label{eq:CS12025-cdf_range_probability}
\end{align}
\solution
The Cumulative Distribution Function (CDF) $F_X(x)$ accumulates the probability density from the lower support limit up to $x$. The probability that a continuous random variable $X$ falls within the interval $[a, b]$ is obtained by evaluating the definite integral of the probability density function $P(x)$, which equals the difference in CDF values:
\begin{align}
P(a \le X \le b) &= \int_a^b P(x) \, dx \\
                 &= \int_{-\infty}^b P(x) \, dx - \int_{-\infty}^a P(x) \, dx \\
                 &= F_X(b) - F_X(a)
\end{align}

% 16. Used in Q28 (Double hashing)
\item Open Addressing Double Hashing Probe Sequence:
\begin{align}
h(k, i) = \left[ h_1(k) + i \cdot h_2(k) \right] \bmod m \label{eq:CS12025-double_hashing}
\end{align}
\solution
In open addressing with double hashing, when a collision occurs at the primary hash location $h_1(k)$, subsequent probe locations are determined by adding a step size governed by the second auxiliary hash function $h_2(k)$ multiplied by the probe index $i$:
\begin{align}
h(k, i) = \left( h_1(k) + i \cdot h_2(k) \right) \bmod m
\end{align}
where $m$ is the hash table size and $i \in \{0, 1, \dots, m-1\}$. To ensure all slots can be probed, $h_2(k)$ must be relatively prime to $m$.

\end{enumerate}
